\section{Hot topics and the Linz school}

Examiners like to end with ``what is new in the field?'' or ``what do you know about our
work?''. This chapter gives you short, safe answers. It covers the Institute for Machine Learning
(IML) at JKU first, because Prof.\ Hochreiter and Prof.\ Klambauer will enjoy talking about their
own work, and then the general hot topics of 2025--2026.

\subsection{The Linz school (IML, JKU)}

\begin{summary}{Linz in two minutes: one timeline}
\begin{tabularx}{\linewidth}{@{}l>{\raggedright\arraybackslash}X@{}}
1991 & Vanishing gradient analysed (Hochreiter, diploma thesis) \\
1997 & LSTM (Hochreiter \& Schmidhuber); Flat Minimum Search \\
2014 & DeepTox wins the Tox21 challenge (Mayr, Klambauer, Unterthiner, Hochreiter) \\
2015/17 & ELU (Clevert et al.) and SELU, self-normalizing networks (Klambauer et al.) \\
2017 & FID, the standard score for generative image models (Heusel et al.) \\
2018/19 & Fr\'echet ChemNet Distance (Preuer et al.); RUDDER, reward redistribution (Arjona-Medina et al.) \\
2020 & Hopfield Networks is All You Need (Ramsauer et al.): attention = Hopfield update \\
2024 & xLSTM (Beck et al.); Vision-LSTM; NeuralDEM with Prof.\ Pirker's group \\
2025 & xLSTM 7B language model and TiRex time-series model (NXAI, Hochreiter's company)
\end{tabularx}
\medskip

\textsf{The common thread:} memory and gradient flow (LSTM, xLSTM, Hopfield), stable deep training
(SELU, flat minima), and AI for science (DeepTox, molecules, physics simulation).
\end{summary}

\pic{Linz is the ``memory school'' of deep learning. Its big ideas are all about keeping
information alive: the error in time (LSTM), the signal in depth (SELU), patterns in memory
(Hopfield), and now long sequences at linear cost (xLSTM).}

\begin{qabox}[Hochreiter, Klambauer]{\pred What do you know about the research of our institute?}
\from{IML research overview · predicted}
\ans The institute is known for LSTM and the analysis of vanishing gradients, for self-normalizing
networks with SELU, and for modern Hopfield networks, which showed that attention is an associative
memory. Recently xLSTM brought LSTM back at scale, with exponential gating and a matrix memory, and
the company NXAI built a 7-billion-parameter xLSTM language model and the TiRex forecasting model.
There is also strong work on AI for science: drug discovery since DeepTox, and physics simulation.
\hook{For me the physics side is the closest: NeuralDEM was done together with my supervisor's group.}
\end{qabox}

\begin{qabox}[Pirker, Klambauer]{\pred What is NeuralDEM, and how does it relate to your thesis?}
\from{NeuralDEM (Alkin, Brandstetter, Pirker et al.~2024) · predicted}
\ans NeuralDEM (Alkin, Brandstetter, Pirker, Lichtenegger and others, 2024) replaces the discrete
element method for industrial particulate flows with a deep-learning surrogate that runs in real
time. It treats the particles as a continuous field and learns their evolution with a
transformer-type model. It is the learned cousin of rCFD: both reuse expensive simulation data to
run fast, rCFD by storing operators, NeuralDEM by learning them.
\follow ``Which is better?'' rCFD conserves the scalar by construction and needs no training;
a learned model can generalize to new states if it has seen enough of them. My thesis shows the
limit they share: a new flow regime is out of the data.
\end{qabox}

\begin{qabox}[Hochreiter]{\pred Why would anyone use an LSTM today, when Transformers exist?}
\from{xLSTM (Beck et al.~2024); LSTM \& Transformers (Hochreiter) · predicted}
\ans Transformers cost $O(N^2)$ in the context length and their memory grows with it at
inference. xLSTM keeps a fixed-size state, so inference is linear in length and constant in memory,
which matters for long contexts, robotics and edge devices. With exponential gating and the mLSTM
matrix memory it reaches Transformer-level language modelling, and mLSTM trains in parallel.
\formula $C_t=f_tC_{t-1}+i_t\mathbf{v}_t\mathbf{k}_t^\top$,\quad
$\mathbf{h}_t=\mathbf{o}_t\odot\dfrac{C_t\mathbf{q}_t}{\max(|\mathbf{n}_t^\top\mathbf{q}_t|,1)}$.
\hook{The same idea works for images (Vision-LSTM) and for time series (TiRex).}
\end{qabox}

\begin{qabox}[Hochreiter]{\pred What is RUDDER?}
\from{RUDDER (Arjona-Medina et al.~2019); Deep RL · predicted}
\ans In reinforcement learning with delayed rewards, credit assignment is slow. RUDDER trains an
LSTM to predict the return of an episode and then redistributes the reward to the steps where the
prediction jumps, i.e.\ to the actions that actually mattered. The expected future reward becomes
close to zero, so learning is much faster.
\pic{A football goal is credited not only to the scorer, but to the pass that made the difference.}
\end{qabox}

\begin{qabox}[Klambauer]{\pred What is the FID?}
\from{FID (Heusel et al.~2017) · predicted}
\ans The Fr\'echet Inception Distance measures how close generated images are to real ones: fit
a Gaussian to Inception-network features of both sets and compute the Fr\'echet distance.
\formula $\mathrm{FID}=\lVert\bm\mu_r-\bm\mu_g\rVert^2+\mathrm{Tr}\bigl(\Sigma_r+\Sigma_g-2(\Sigma_r\Sigma_g)^{1/2}\bigr)$.
\follow The same idea for molecules is the Fr\'echet ChemNet Distance.
\end{qabox}

\subsection{Hot topics 2025--2026}

\begin{summary}{Six trends to name}
\begin{enumerate}
\item \textsf{Reasoning LLMs:} models trained with reinforcement learning on verifiable rewards
  (maths, code) to think longer before answering: test-time compute.
\item \textsf{Alignment:} RLHF and its simpler successor DPO learn from human preferences.
\item \textsf{Efficient sequence models:} xLSTM, Mamba (state space), linear attention; and
  mixture-of-experts Transformers that activate only a few experts per token.
\item \textsf{Generative models:} diffusion and flow matching for images, video, molecules.
\item \textsf{Agents and retrieval:} LLMs that call tools and read documents (RAG).
\item \textsf{AI for science:} protein structure (AlphaFold 3), weather (GraphCast, GenCast),
  neural surrogates for CFD and particle flows (NeuralDEM), foundation models for physics.
\end{enumerate}
Plus regulation: the EU AI Act (risk-based rules), which links to your AI and Law courses.
\end{summary}

\pic{The field is moving from ``bigger models'' to ``smarter use of compute'': think longer
(reasoning), use only the needed part (mixture of experts), remember cheaply (xLSTM, Mamba), and
ask tools instead of guessing (agents).}

\begin{qabox}[general]{\pred What is flow matching? (You had it in ML: Advanced Techniques.)}
\from{ML: Advanced Techniques · flow matching unit · predicted}
\ans A generative model that learns a velocity field which moves noise to data along a path. With
the straight path $\mathbf{x}_t=(1-t)\mathbf{x}_0+t\mathbf{x}_1$ the target velocity is simply
$\mathbf{x}_1-\mathbf{x}_0$, so training is a plain regression. Sampling solves the ODE
$\mathrm{d}\mathbf{x}/\mathrm{d}t=\mathbf{v}_\theta(\mathbf{x},t)$.
\formula $L(\theta)=\mathbb{E}_{t,\mathbf{x}_0,\mathbf{x}_1}\bigl\lVert\mathbf{v}_\theta(\mathbf{x}_t,t)-(\mathbf{x}_1-\mathbf{x}_0)\bigr\rVert^2$.
\hook{Diffusion models are the same family with a noisy, curved path.}
\end{qabox}

\begin{qabox}[general]{\pred Diffusion models in two sentences?}
\from{ML: Advanced Techniques · diffusion (VDM) · predicted}
\ans Add Gaussian noise step by step until the data is pure noise, and train a network to predict
the added noise; sampling runs the process backwards. Loss:
$\mathbb{E}\lVert\bm\epsilon-\bm\epsilon_\theta(\sqrt{\bar\alpha_t}\mathbf{x}_0+\sqrt{1-\bar\alpha_t}\bm\epsilon,t)\rVert^2$.
\end{qabox}

\begin{qabox}[general]{\pred How are LLMs aligned? RLHF versus DPO.}
\from{Current literature: RLHF, DPO · predicted}
\ans RLHF: train a reward model on human preference pairs, then optimize the LLM against it with
RL (PPO), with a KL penalty to stay close to the original model. DPO skips the reward model and
optimizes the preferences directly with a classification-like loss.
\formula DPO: $L=-\log\sigma\Bigl(\beta\log\dfrac{\pi_\theta(y_w|x)}{\pi_{\mathrm{ref}}(y_w|x)}-\beta\log\dfrac{\pi_\theta(y_l|x)}{\pi_{\mathrm{ref}}(y_l|x)}\Bigr)$.
\hook{Reasoning models go one step further: they learn from rewards that can be checked automatically.}
\end{qabox}

\begin{qabox}[general]{\pred What are reasoning models and test-time compute?}
\from{Current literature: reasoning models · predicted}
\ans Models that write out intermediate steps (chain of thought) before the answer, trained with
reinforcement learning on tasks whose answers can be verified, like maths or code. Spending more
computation at inference, by thinking longer or sampling several solutions, improves accuracy:
test-time compute becomes a new scaling axis next to model size and data.
\end{qabox}

\begin{qabox}[general]{\pred Mixture of experts and state space models?}
\from{Current literature: MoE, Mamba · predicted}
\ans Mixture of experts: many feed-forward ``experts'', a router sends each token to the top-$k$,
so the model has many parameters but uses few per token:
$\mathbf{y}=\sum_{i\in\text{top-}k}g_i(\mathbf{x})E_i(\mathbf{x})$. State space models like Mamba:
a linear recurrence $\mathbf{h}_t=\bar A\mathbf{h}_{t-1}+\bar B\mathbf{x}_t$ with input-dependent
(selective) parameters; linear in sequence length, like xLSTM.
\end{qabox}

\begin{qabox}[Pirker, general]{\pred Where is AI for simulation going?}
\from{AI for science literature; NeuralDEM · predicted}
\ans From single surrogates to foundation models for physics: neural operators and transformer
surrogates trained on many simulations, weather models that beat numerical forecasts in speed,
and industrial real-time twins like NeuralDEM. The open problems are the ones my thesis met:
conservation, generalization to new regimes, and trust, i.e.\ knowing when the surrogate is out of
its data.
\end{qabox}
